%!TEX root = ../template.tex
%%%%%%%%%%%%%%%%%%%%%%%%%%%%%%%%%%%%%%%%%%%%%%%%%%%%%%%%%%%%%%%%%%%%
%% abstract-en.tex
%% NOVA thesis document file
%%
%% Abstract in English([^%]*)
%%%%%%%%%%%%%%%%%%%%%%%%%%%%%%%%%%%%%%%%%%%%%%%%%%%%%%%%%%%%%%%%%%%%

\typeout{NT FILE abstract-en.tex}%

Formal Languages and Automata Theory (FLAT) is rigorous and often perceived as abstract, which complicates early learning. This dissertation extends 
the \textbf{OCamlFLAT} library with support for \emph{attribute grammars}, including synthesized and inherited attributes, evaluation strategies, and 
static checks to detect common specification errors. The contribution focuses on a reliable, executable core and clear interfaces.

The original plan included enhancements to the \textbf{OFLAT} web application and a pedagogical evaluation. In the final scope, the work is limited to 
the extension of the \textbf{OCamlFLAT} library with the necessary structures needed for later integration. The graphical components and the empirical study of pedagogical 
impact are \emph{out of scope} and are outlined as future work.

The dissertation reviews the relevant background on attribute grammars within the theory of computation, surveys related educational tools, 
and details the architecture and implementation of the OCamlFLAT extensions. By stabilising the core and defining integration points, the work 
provides a basis for subsequent visualisation and classroom use.


% Palavras-chave do resumo em Inglês
% \begin{keywords}
% Keyword 1, Keyword 2, Keyword 3, Keyword 4, Keyword 5, Keyword 6, Keyword 7, Keyword 8, Keyword 9
% \end{keywords}
\keywords{
  Formal Languages \and
  Automata Theory \and
  OCaml \and
  OFLAT \and
  OCamlFLAT \and
  Functional Programming \and
  Formal Grammars
}
